Distributed Key Generation (DKG) schemes underlie many
multi-party cryptographic primitives in use today,
such as threshold signature schemes~\cite{KomloG20,BellareCKMTZ22,Lindell22} and distributed randomness beacons~\cite{le19}.
A DKG allows a set of $n$ participants to cooperate to generate a keypair,
such that the public key represents the entire set of participants,
and each participant holds a share of the corresponding secret key.
DKGs ensure that the public key is agreed upon by all (honest) parties at the end of the protocol,
but no single party knows the corresponding secret key.
However, this secret key can be recovered given a threshold $t$ number of cooperating parties,
where $t \leq n$.

Despite widespread interest in DKGs,
existing constructions fall short in either their complexity or composability.
Most glaringly is Pedersen-DKG~\cite{Pedersen91} and variants of Pedersen DKG
such as PedPop~\cite{KomloG20,BellareCKMTZ22}.
While Pedersen-DKG is appealing because of its simplicity and efficiency,
its security to date has only been proven in the context of specific applications~\cite{GennaroJK02,GurkanJMMST21,BellareCKMTZ22},
as it cannot be proven secure in a standalone fashion~\cite{GennaroJKR07}.
More specifically,
it is possible that the use of Pedersen-DKG may be secure in some settings, but not in others.
As such, the risk exists that Pedersen-DKG might be misused in contexts for which it cannot been proven secure.
While many alternative DKG constructions have been presented in the literature~\cite{GennaroJKR07,Groth21,GurkanJMMST21,BonehS2020,NejiBR16,Lindell22,DasYXMKR21,BachoL22,AbrahamJMMS2022},
these constructions incur undesirable protocol complexity or require pairings.

Conversely, growing interest in threshold schemes points
towards the need for DKGs that are friendly to standardization efforts~\cite{NIST22,NIST23} and
securely composable across a range of use cases.
For example, the FROST Schnorr threshold signature scheme~\cite{KomloG20,BellareCKMTZ22},
for which an IETF draft exists~\cite{Connolly22},
is written in a way to be agnostic to any particular key generation scheme.
Ideally, a candidate DKG to be standardized should be applicable not only to
a specific threshold signature scheme,
but also to other applications that require secret key material to be distributed
among a set of trusted parties~\cite{DavidsonSSW22}.

\paragraph{Proving the Security of DKGs.}
However,
challenges arise when proving the security of a DKG in a standalone manner;
i.e., that the DKG is secure in \emph{any} setting in place of its target (single-party) key generation scheme.
Although existing simulatability-based notions for DKGs exist~\cite{GennaroJKR07,GurkanJMMST21,Lindell22,BachoL22,BonehS2020,BendlinKP13},
these notions prove to be either incomplete or insufficient.
For example,
many simulatability-based notions assume important details such as the
consistency of honest participants' state or that adversarial players do not cause the protocol to abort.
However, such assumptions are \emph{not} guaranteed in a distributed, multi-party setting.
Notably, participants that force the protocol to abort during the last round
can be employed as a vector for a key-bias attack~\cite{GennaroJKR07}.
Under the hood, many definitions assume the protocol to be robust;
i.e., that the protocol will always complete successfully in spite of misbehaving parties.
However, doing so increases the complexity for the protocol
by requiring additional complaint and voting rounds, which are difficult to
implement correctly in practice.


\subsection{Our Contributions}

\paragraph{A Composable and Generic DKG.}
In this work, we present a novel and simple generic DKG construction that is a good candidate for ongoing threshold standardisation efforts.
Our construction relies on three simple cryptographic building blocks:
a secure hash function,
an \emph{Aggregatable Verifiable Secret Sharing} (AgVSS) scheme,
and a \emph{Non-Interactive Key Exchange} (NIKE) scheme.
The ability to aggregate Feldman's VSS~\cite{Feldman87}  has been extensively employed in prior DKG
constructions~\cite{Pedersen91,GennaroJKR07,KomloG20,GurkanJMMST21,Lindell22};
we simply formalize the notion of a VSS that is aggregatable,
and present new security notions to reflect this property.
We additionally provide formal definitions and notions of security for a NIKE,
as well as concrete AgVSS and NIKE constructions.

\begin{table}[t]
  \centering
  \small
  \begin{tabular}{c|c|c|c|c|c|c}
	\hline
    & \multicolumn{2}{c|}{Composable \& VSS}
    & \multicolumn{2}{c|}{Composable \& PVSS} &
    Not Composable \\
	\hline
    & \GJKRConstruction~\cite{GennaroJKR07}
    &  \RG
    & X
    & X
    &  $\pedpop$~\cite{KomloG20,BellareCKMTZ22}
      \\
	\hline
    \textbf{Optimistic Rounds} & 3 & 3  & X & X & 2 \\
    \textbf{Total Rounds} & 6 & 4 & X & X & 2 \\
    \textbf{Robustness} & Strong & Weak &  X & X &No \\
    \textbf{Bandwidth} & $2nt + 3n$ &  $nt+ 5n$ & X & X & $nt +3n$  \\
    \textbf{Computation } & $2nt + 2n + t$ &  $nt + 2n$ & X & X & $nt + 2n$   \\
    \textbf{Assumption} & DLP & CDH & X & X & KoE/AGM   \\
    \textbf{Threshold} & $n \geq 2t-1$ &  $n \geq 2t-1$ & X & X & $n \geq t$  \\
\hline
	\end{tabular}
  \smallskip
  \caption{Comparison between \RG and related synchronous constructions that are secure in the random oracle model,
  that are based on discrete logarithm assumptions,
  do not require pairings,
  and have formal proofs of security.
  Composable denotes if the construction has a standalone proof of security
  and can be used in any context that its target (single-party) key generation
  algorithm is used.
  %as opposed to requiring a (new) proof of security for each use case.
  Optimistic rounds considers only when no participant cheats;
  total rounds includes additional rounds required to recover from players that
  misbehave during the protocol.
  Robustness indicates whether the protocol always succeeds assuming some
  number of honest players (strong) or allows non-fatal aborts (weak).
  Bandwidth efficiency denotes the number of group and field elements sent and received by each participant throughout the protocol.
  Computational efficiency denotes the number of exponentiations required (i.e., scalar multiplications in an elliptic curve group).
  $n$ is the total number of participants in the protocol,
  $t$ is the threshold number required to recover the secret key,
  and the number of corrupted parties is assumed to be \emph{at most} $t-1$.
  $k$ is security parameters for the respective scheme~\cite{DasYXMKR21}.
  %$k,r$ are security parameters for their respective schemes~\cite{DasYXMKR21,Lindell22,Fischlin05}.
  DLP is the Discrete Logarithm Problem;
  KoE is the Knowledge of Exponent model; AGM is the Algebraic Group Model;
  CDH is the Computational Diffie-Hellman Problem.
  }
  \label{tbl:intro-comp}
\end{table}



\medskip

\noindent \textit{Why a Generic Construction?}
The goal for presenting a generic DKG as opposed to a monolithic construction is to enable flexibility for implementations and future standardization efforts.
For example, while we define concrete AgVSS and NIKE constructions,
future implementations may instead wish to employ a standardized
NIKE~\cite{hpke} instead.
Further, as NIST is currently in the process of standardizing threshold primitives~\cite{NIST23},
it may be worthwhile to leverage future VSS standards, which may possibly diverge from the concrete VSS defined in this work.
Or, implementations may wish to define a construction that employs publicly verifiable randomness,
a polynomial commitment scheme that is more efficient,
or derives secret values deterministically, such as from a seed phrase.
The utility of a generic DKG is that so long as the chosen building blocks are secure,
any instantiation of our generic DKG using these building blocks will remain secure.

\paragraph{Concrete Instantiation: \RG.}
We then describe a concrete instantiation of our generic construction that we call
\RG (\underline{S}ynchronous, dis\underline{T}ributed, and
\underline{O}ptimized gene\underline{R}ation of key \underline{M}aterial).

The security of \RG reduces to the Computational Diffie-Hellman (CDH) problem in the random oracle model.
\RG is securely composable with existing threshold cryptosystems such
as threshold Schnorr signatures~\cite{KomloG20,Connolly22} and similar
discrete-log based systems.
We compare in Table~\ref{tbl:intro-comp} the practicality and security of \RG compared to related constructions in the literature that require only standard (i.e., non-pairing based) groups.
\RG is simpler and therefore (slightly) more efficient than related syncronous DKGs in the discrete-logarithm setting with similar
security guarantees~\cite{GennaroJK02}.
Its efficiency improvement is due to the use of a NIKE,
which allows the parties to derive
a ``tweaked'' key at the end of the protocol, preventing adaptive attacks.





\paragraph{Foramlized Notions of Security for DKGs.}
Our goal is to prove the security of our generic construction in such a way that
(i) demonstrates the DKG can be securely used in any setting where its target key generation scheme is used,
and (ii) captures implicit notions such as handling adversarial aborts and honest participant consistency.
Towards this end,
we introduce formal game-based security notions for DKGs
that generalize and extend prior notions in the literature.
Our games make explicit important details such as communication patterns between honest and adversarial players,
where prior notions have omitted some details.
For example,
our games explicitly model the adversary as \emph{rushing},
in that it can always query for honest player's contributions before producing its own.
We then prove the security of our generic DKG with respect to these notions.

For a DKG to be secure, we require that the DKG  be strongly or weakly robust,
zero-knowledge, and
indistinguishable.
Strong versus weak robustness captures the distinction between a protocol that always succeeds assuming some threshold of honest participants (strong robustness),
and one which allows for non-fatal aborts that do not impact the security of the scheme (weak robustness).
Similar to prior simulation-based notions,
zero-knowledge requires that the DKG is simulatable with respect to a public
key.
Indistinguishability requires that the output of a DKG be indistinguishable from its target (single-party) key generation scheme.
By fulfilling the notion of indistinguishability,
the DKG can be securely employed in any setting where the target (single-party) key generation scheme could be used.
However,
for cryptosystems that ensure secret and public keys have a unique correspondence (such as discrete-log based schemes),
zero-knowledge implies indistinguishability.
%

These notions could be captured by an alternative model such as Universal Composability~\cite{Canetti01}.
We instead opt for a game-based approach,
so we can explicitly and independently define each notion.
%However, because these notions capture
%perfect simulatability with respect to a target key generation algorithm,
%DKGs that achieve our notions of security without rewinding achieve UC security~\cite{KushilevitzLR10}.



\paragraph{What We Do Not Do.}
Our definitions and constructions assume a strictly synchronous setting;
i.e., that every participant in the protocol terminates a round before beginning the next.
While asynchronous DKGs
exist~\cite{NejiBR16,DasYXMKR21,AbrahamJMMST21,Kokoris-KogiasM20,KateHG12,YurekLFKM22},
these constructions require the tradeoff of increased complexity.
For similar reasons, we require the existence of a secure broadcast channel.
%While alternative schemes do not~\cite{YurekLFKM22},
%the tradeoff is higher complexity and decreased modularity.
%We target the setting where key generation is performed infrequently or in a more controlled environment.

Additionally, while some prior DKGs build upon a publicly verifiable secret
sharing scheme (PVSS)~\cite{Groth21,BonehS2020},
we explicitly do not,
due to the tradeoff in increased implementation complexity.
Furthermore,
we do not consider the large-scale
setting~\cite{CannyS04,ZhangXHSZ22,YurekLFKM22}, although such optimizations
could be employed to define a more efficient instantiation of our
generic construction.

We leave proofs of adaptive security for our generic construction for future work.
Techniques to achieve tight proofs of adaptive security in stronger security
models (such as the AGM) have recently been demonstrated in concurrent work~\cite{BachoL22,AbrahamJMMS2022};
we expect these techniques to similarly apply to our constructions.

%Finally, our proof for our generic construction is in the ROM.
%However, we believe recent results lifing some proofs to the QROM~\cite{BonehDFLSZ11,GriloHHM21} to
%likely be applicable to our proof as well.


\paragraph{Open Questions.}
Our generic construction assumes a one-to-one correspondence between public and private keys.
While \RG is secure assuming the hardness of the Computational Diffie-Hellman (CDH) problem,
it is possible our generic construction can be instantiated by quantum-secure primitives that define a homomorphism between the secret and public domains.
For example, performing rounds in a sequential manner may allow our generic construction to be instantiated by CSIDH-based primitives~\cite{CastryckLMPR18,BeullensDPV21}.
However, because of the difficulty of instantiating a NIKE using lattice-based primitives,
it is an open question as to how our generic construction can be extended to lattice-based primitives.


\paragraph{Summary of Our Contributions.}
In this work, we present the following contributions.

\begin{itemize}
  \item We introduce the notion of an aggregated Verifiable Secret Sharing (AgVSS) scheme,
    building upon its implicit use in prior literature.
    We then give concrete AgVSS and Non-Interactive Key Exchange (NIKE) constructions.
  \item We present formal game-based definitions for the security of a DKG,
    that makes explicit details such as communication patterns between honest and corrupted players,
    expectations of participant consistency and requirements of key validity.
    Further, our separation of the notion of robustness into a strong and weak variant allows for
    a simpler and more efficient DKG construction.
  \item We give a generic DKG construction,
    and prove its security using our new notions.
    Our generic construction requires a secure hash function, an AgVSS, and a NIKE.
  \item We introduce \RG, a concrete instantiations of our generic
    construction.
    The security of \RG reduces to the Computational Diffie-Hellman (CDH)
    problem in the random oracle model.
\end{itemize}

